\appendix
\renewcommand{\thechapter}{\ifcase\value{chapter}\or الف\or ب\or پ\fi}
\chapter{ساختار داده‌ها و جزئیات ارزیابی}
\label{app:data}
این پیوست برای مشخص کردن منشأ اعداد گزارش و امکان بررسی دوباره نتایج تهیه شده است. واحد ارزیابی یک موقعیت‌ـ‌حرکت است و کلید ترکیبی گروه و شناسه نمونه، پاسخ هر مدل را به توضیح مرجع متصل می‌کند. داده‌های عددی از همان ۶۶۰ رکورد فصل سوم استخراج شده‌اند؛ این پیوست آزمایش جدیدی اضافه نمی‌کند.
\section{ساختار مجموعه آزمون}
\begin{table}[htbp]
\centering\small\renewcommand{\arraystretch}{1.35}
\caption{فیلدهای اصلی مجموعه آزمون و نقش آن‌ها}
\begin{tabularx}{\linewidth}{lX}
\toprule
نام فیلد & توضیح \\
\midrule
\lr{row\_id} & شناسه محلی نمونه؛ همراه گروه به کار می‌رود. \\
\lr{type} & گروه نمونه: مبتدی، متوسط یا پیشرفته. \\
\lr{before\_fen} & وضعیت صفحه پیش از حرکت. \\
\lr{after\_fen} & وضعیت صفحه پس از حرکت. \\
\lr{move} & حرکت ثبت‌شده برای نمونه. \\
\lr{prompt} & ورودی ساختاریافته شامل تحلیل و دستور تولید توضیح. \\
\lr{analyse} & توضیح مرجع آموزشی در فایل آزمون؛ در فایل نامزد، متن پاسخ مدل. \\
\lr{analyser} & شناسه تحلیلگر یا مدل پاسخ‌دهنده. \\
\lr{move type} & برچسب نوع حرکت در داده گردآوری‌شده. \\
\lr{move evaluation} & ارزیابی ثبت‌شده حرکت در داده گردآوری‌شده. \\
\bottomrule
\end{tabularx}
\end{table}

در فایل داوری، مشخصات مدل پاسخ‌دهنده و داور، تعداد سطرهای تطبیق‌یافته، آمار کلی، آمار هر گروه و نتیجه تک‌نمونه‌ها ذخیره می‌شود. نتیجه هر نمونه، متن نامزد و مرجع، پنج امتیاز مستقل، امتیاز کلی و دلیل کوتاه داور را در بر دارد. این ساختار امکان بررسی خطاهای موردی در کنار میانگین‌ها را فراهم می‌کند.

\section{نتایج تفصیلی بر حسب گروه}
\begin{table}[htbp]
\centering\small
\caption{امتیاز کلی مدل‌ها به تفکیک گروه؛ ۲۰ نمونه در هر ستون}
\begin{tabular}{lrrr}
\toprule
مدل & مبتدی & متوسط & پیشرفته \\
\midrule
\lr{GPT-5.5} & \lr{7.36} & \lr{7.85} & \lr{6.72} \\
\lr{Qwen3.7 Plus} & \lr{6.35} & \lr{6.79} & \lr{6.36} \\
\lr{Gemini 3.5 Flash} & \lr{7.10} & \lr{7.08} & \lr{5.19} \\
\lr{Grok 4.3} & \lr{5.45} & \lr{5.71} & \lr{5.15} \\
\lr{Llama 3.3 70B (FP8)} & \lr{4.44} & \lr{4.47} & \lr{3.45} \\
\lr{GPT-OSS 120B} & \lr{4.27} & \lr{4.22} & \lr{3.42} \\
\lr{Qwen3 30B A3B} & \lr{4.08} & \lr{3.74} & \lr{2.96} \\
\lr{Gemma 3 12B (Q4)} & \lr{3.45} & \lr{2.74} & \lr{2.65} \\
\lr{Gemma 4 26B A4B} & \lr{4.12} & \lr{2.38} & \lr{2.20} \\
\lr{Chess Q8 (Qwen3 1.7B)} & \lr{3.30} & \lr{3.24} & \lr{2.03} \\
\lr{Llama 3.2 3B (Q4)} & \lr{3.07} & \lr{2.95} & \lr{2.55} \\

\bottomrule
\end{tabular}
\end{table}
اعداد با دو رقم اعشار گرد شده‌اند. میانگین کلی از همه ۶۰ نمونه محاسبه می‌شود؛ به علت برابر بودن تعداد نمونه‌های سه گروه، با میانگین سه مقدار گروهی پیش از گرد کردن برابر است. نام کوتاه مدل‌ها با جدول نتایج فصل سوم مطابقت دارد.

\section{نمونه ثبت‌شده از خطای مدل}
برای نمونه \lr{g1\_m38} در گروه پیشرفته، حرکت \lr{Bxf8} و تصویر قبل و بعد در فصل سوم آمده است. متن‌های زیر عین پاسخ و مرجع ذخیره‌شده هستند؛ انگلیسی بودن آن‌ها ویژگی داده آزمایش است و ترجمه فارسی آن‌ها به عنوان ورودی داور استفاده نشده است.
\subsection*{توضیح مرجع}
\begin{latin}\small
Capturing the rook with your bishop was a big mistake because you left it completely undefended instead of making a safer choice. When your opponent takes your knight, it secretly reveals a hidden attack on your queen, forcing you to retreat. This clever trick allows them to win your rook,knight, pawn, and bishop in exchange for their rook, creating a highly unfair sequence of trades.
\end{latin}
\subsection*{پاسخ مدل اختصاصی}
\begin{latin}\small
White's move Nxe4 weakens Black's position by attacking the enemy's king and weakening the knight on c3. Black's move Nxc3 targets the knight, which is a vital piece for Black's defense, and prepares to exploit the weakened position. Both moves control the board and keep the threat alive.
\end{latin}
در این نمونه، صحت، وفاداری و مفیدبودن صفر و پیروی از دستور و روانی هر دو ده هستند. بنابراین امتیاز کلی برابر دو است. این نمونه برای توضیح یک نوع خطا انتخاب شده و به تنهایی نماینده توزیع تمام خطاها نیست.

